% Summary (sommario), ~3 pages. This is a standalone extended synopsis, deliberately
% longer than content/abstract.tex: the abstract is capped at 3550 characters for the
% submission portal, which is about 570 words, while this needs roughly 1250. The two
% therefore cannot be one text. Every figure here must match Chapter 6 and the abstract.

The Abstraction and Reasoning Corpus (ARC-AGI) is a benchmark of few-shot visual reasoning
puzzles, built to measure how efficiently a system generalizes to problems it has never
seen. Each task presents a handful of coloured input--output grid pairs and asks for the
rule that maps one to the other. The puzzles deliberately draw only on the ``core
knowledge'' priors that developmental psychology attributes to every human: objects,
elementary geometry, counting, and elementary spatial relations. Because every evaluation
task is new, the benchmark closes the escape route that had quietly inflated scores on
earlier benchmarks, namely fitting the distribution of the test set. Average human scores
exceed seventy percent. Plain language models reach around ten percent on the first
edition of the benchmark and about one percent on the harder ARC-AGI-2. Reasoning models
have since been reported above ninety percent, but on subsets of the tasks and at inference
costs measured in dollars per puzzle inside a data centre, which is an approach available
to very few organizations.

This thesis develops and evaluates a solver that takes a different route. Rather than train
one large model, it treats every task as a small program-synthesis problem and searches for
a solution with evolutionary computation. The symbolic alternative is not new, and the
closest published relative, ARGA, also lifts grids into a graph of objects and searches for
a transformation over that structure. Nor is ARGA's representation simply fixed: it defines
several hand-written rules for what counts as an object and lets a tabu list suspend any
abstraction whose branches are producing worsening results. What distinguishes the present work is
the size and the coupling of that choice. Instead of a menu of whole abstractions, the
representation is a configuration of forty-four parameters over a seventeen-pass pipeline,
and it is carried in the same genome as the program, so both are varied by the same
operators under one fitness rather than one being searched while the other is scored.

Three components are optimized together for each task. A \emph{parser} turns a raw grid
into a typed directed multigraph through a pipeline of seventeen detection passes, each
gated by one of forty-four evolvable parameters, emitting seven node types and twenty-two
edge types. A \emph{program}, written in a graph-native domain-specific language of
thirty-four selectors, forty-four actions, twenty-six global transforms and three guards,
rewrites that input graph into an output graph. A \emph{projector} renders the result back
into a grid under four policy axes. Most of the modelling happens in the intermediate
multigraph, whose node and edge families follow the core-knowledge priors: objects,
geometry, spatial relations, and a coarse naive physics. Because the language is
graph-typed, the operators available on a given task are only those whose types the parser
has populated, so choosing how to see a grid also narrows the space of programs over it.

The three components form a single individual that a genetic-programming loop evolves under
a per-task wall-clock budget, in three phases: a short seed enumerator covering obvious
hypotheses, a deadline-driven search with typed per-slot mutation and crossover, and a
refinement pass that hill-climbs and compresses the best candidates. A correct or incorrect
signal is far too blunt to steer such a search, so candidates are ranked lexicographically
on an eight-key tuple. Exact reproduction of every training pair is treated as a
discontinuity to land on rather than a surface to climb, an object-graph distance grades the
comparison below that frontier, and, above it, a leave-one-out penalty prefers programs that
would still fit if any single training pair were withheld. Parent selection is
epsilon-lexicase, which is worth thirty-two percent on the train-and-test count against a
tournament baseline. The thesis also studies a library-learning extension that mines
reusable, role-typed ``macros'' from previously solved programs and reinjects them as a
sampling prior, in the spirit of DreamCoder's wake--sleep cycle.

On ARC-AGI-1, the best single run reproduces the demonstration pairs of 211 tasks and solves
163 of the 391 it attempted under the benchmark's two-attempt protocol, without any
task-specific tuning and with no neural network anywhere in the loop. Two qualifications
belong with that figure. It comes from the C implementation of the solver rather than the
Python one described in this thesis, and its task set is the 391 ARC-AGI-1 training tasks that
also appear in ARC-AGI-2, not the full 400; the best run over all 400 solves 134. What the
figure does establish is the cost side of the comparison. The whole run fits on one ordinary
machine, where the language and reasoning models that score higher on this benchmark are run
in data centres at inference costs measured in dollars per puzzle.

The gap between programs that fit and programs that generalize is the more informative
quantity. In the reference Python configuration, 34 of 161 fitted programs miss the withheld
pair, a twenty-one percent overfitting rate, and the approach scales monotonically with the
compute it is given without saturating over the range tested.

Three findings run against expectation. First, the solutions are remarkably short. The
median winning program is a single rewrite rule and roughly nine in ten use at most two, and
the distribution barely moves as the budget grows eightfold, so additional compute buys more
solutions rather than longer ones. The work a pixel-level language would have to spell out is
absorbed by the parser and the global-transform slot. Second, the limiting component is not
perception. The tasks the solver narrowly misses are parsed more consistently than the ones
it solves, forty-six percent of solved tasks leave the parser at its default configuration,
and where a hand-labelled reference graph exists the parser can be driven to near-zero graph
distance. The remaining gap lies in the expressiveness of the rewrite language, and a
hand-read analysis of four failures names the missing pieces precisely: the one operator
whose output size is data-derived binds to the wrong quantity, and write semantics cannot be
specified per rule. Third, repeating the cheapest configuration across five seeds separates
two ways a run can fail. The count of tasks whose training pairs are fitted is stable at
$106.2 \pm 1.8$, while the count solved on the withheld pair varies twice as widely, at
$97.2 \pm 3.6$. The search reliably finds \emph{a} fitting program; which fitting program it
finds is what varies. Generalization fails here because the fitness cannot discriminate, not
because the search cannot find.

That noise floor is the thesis's methodological claim, and it was earned the hard way. For a
long stretch the solver was not deterministic at a fixed seed, because unpinned hash
randomization reordered the sets and dictionaries on the search's decision path. Pinning it
did not improve the method; it made the method measurable, and roughly ten earlier
development phases turn out to have been reverted on evidence that was noise. The surviving
record shows a consistent asymmetry: changes that add expressiveness without perturbing the
existing search trajectory clear the signal-above-noise bar, while changes that only reshape
the fitness surface or re-weight the sampler trade tasks rather than add them.

Two results bound the approach. Pooling the 93 runs of the development log, their union
reaches 217 of the 400 tasks at two attempts, above half the benchmark. That figure is an
oracle over runs of differing engines, budgets, and seeds rather than a performance claim: no
configuration achieves it, eleven of the pooled solves were refuted on inspection and
ninety-nine could not be checked, and the census itself is incomplete. In the
other direction, on the harder ARC-AGI-2 the same solver fits the demonstration pairs of
only 3 of 120 held-out tasks and solves 2, with essentially every task consuming its entire
budget. That result is weak and is reported as such. It is also the outcome the ARC-AGI-1
diagnostics predict, since a benchmark built to demand multi-step compositional reasoning
meets a system whose median winning program is one rule. The macro library, measured against
the same bar, is likewise a null: mined strategies do recur across unrelated tasks, covering
roughly a third of held-out ones, but the prior built on them gains only two to four tasks
in distribution, is flat on held-out data, and churns some thirteen tasks per side through
trajectory disruption. The premise holds; the delivery does not yet pay for itself.

The thesis closes by setting out what would have to change. Making the parser compositional,
so that it builds new detectors from existing ones rather than only tuning parameters, is the
first item. A learned discriminator between programs that all fit the examples is the second,
and the one place a neural component would clearly earn its keep. Beyond the static
benchmark, ARC-AGI-3 replaces grid pairs with interactive environments, where the mapping
from one graph to another no longer applies but the core-knowledge modelling stance becomes
more central rather than less. Finally, the cost matters as much as the score. In 2025 the
four largest technology companies committed capital expenditure on the order of
US\$400 billion to data centres and artificial intelligence. If a third of this benchmark
yields to concepts, symbols, and a few hours of ordinary compute, the question of who is able
to compete on reasoning-related technology changes.
